\section{Source conventions, prior results, and remaining limits}
\label{b-interpretation}

\subsection{Collective completion is not an individual self-energy statement}
Under the Gaussian-domination hypotheses, the exact source conclusion is
\begin{equation}
 E_\nu(q)-U_{\rm nn}
 +\frac\alpha2D_L^{\rm band}(\nu,\nu)
 \ge E^{\vac}+m_*|q|.
 \label{b-collective-completion}
\end{equation}
The band form is a collective source energy, containing nuclear cross
terms as well as individual diagonals. If
$G_L^{\rm band}=V^{-1}\sum_{0<|k|\le2\Lambda_f}v_ke^{ikx}$ and
$U_{\rm nn}^{\rm band}=\alpha\sum_{i<j}Z_iZ_j
G_L^{\rm band}(R_i-R_j)$, then equivalently
\[
 E_\nu(q)+\frac\alpha2\kappa_{\Lambda,L}\sum_iZ_i^2
 \ge E^{\vac}+m_*|q|+U_{\rm nn}-U_{\rm nn}^{\rm band}.
\]
The last difference has no fixed sign. The kernel
$G_L-G_L^{\rm band}$ has zero mean and a positive Coulomb singularity,
so it must be negative somewhere. Positive Fourier coefficients do
not imply pointwise positivity. The heat argument is needed to control
this mismatch for many point nuclei. A statement that the nuclei plus
their individual band self-energies never lower the energy below the
vacuum is valid for one nucleus or for a consistently band-regulated
nuclear repulsion, but is too strong for the original point-source
repulsion. The small-box example~\eqref{eq:small-box} is an explicit
counterexample to its unrestricted version.

A common scalar shift cancels from every relative energy. Adding
$\sum_i\varepsilon_{Z_i}(\Lambda_f)$ instead changes the source
reference. Such a species-dependent calibration must be independently
specified, volume independent, and independent of configurations and
states. It is distinct from subtracting the infinite classical self-energy
of an individual point nucleus. No additional nuclear calibration is
used here, and neither an arbitrary configuration subtraction nor a
finite-box source ground energy is used to force the bound.

\subsection{The exact models in the relevant literature}
Lieb--Yau~\cite{LY1988} prove stability and instability for relativistic
particle kinetic energy and point Coulomb sources; there is no
pair-creating Dirac field or quantized transverse photon field in that
model. Lieb--Loss~\cite{LL2002} include quantized photons and point
nuclei in a no-pair theory using the positive spectral subspace of a
field-dependent Dirac operator. Positrons are excluded; at physical
$\alpha$ their stated range includes $Z\le42$. Lieb's review
\cite{Lieb2002} explains these distinctions. Neither no-pair theorem
is applied to the full Fock Hamiltonian here.

J\"akel--Wreszinski~\cite{JW2018} treat finite-volume, ultraviolet-
regularized Coulomb-gauge fermion and photon fields. Their fixed-regulator
vacuum-energy shift is allowed to depend on volume. It is not the
uniform point-nucleus comparison proved above. The lower-bound square
completion here is an explicitly proved operator inequality; it is
not justified by assuming that shifted current-dependent oscillators
obey canonical commutation relations.

Lieb--Siedentop~\cite{LS2000} already distinguish a positive
density-product interaction from its normal-ordered form and analyze
the resulting quadratic operator with a self-consistent ordering.
Our reference is prescribed to be the free sea, so its exact one-body
and scalar corrections are retained rather than replaced by that
self-consistent construction. Bach--Barbaroux--Helffer--Siedentop
\cite{BBHS1998} prove matter stability for a relativistic
Hartree--Fock electron--positron functional with point nuclei and
classical magnetic fields. Their reference splitting and variational
class differ from this fully correlated quantized-field Hamiltonian.
Hainzl--Lewin--S\'er\'e~\cite{HLS2004,HLS2005,HLS2008} establish
stable polarized vacua and atomic or molecular existence in specified
BDF approximations. Hainzl--Lewin--Solovej~\cite{HLSol2007} construct
a self-consistent no-photon mean-field vacuum and its thermodynamic
limit. These do not identify that vacuum with the free sea in a
different ordering, or supply a theorem for arbitrary correlated states.
The periodic Coulomb--Kato estimate used here is part of that established
toolkit; the explicit constant is proved in Lemma~\ref{lem:pkato}.

The point-nucleus Dirac--Coulomb distinguished-extension theory
\cite{ELS2021} concerns one-particle operator domains and critical
couplings. At our finite fermion band the compressed source matrices
are bounded, so the proof does not invoke a singular one-particle
extension. This regularity benefit is lost on removing the cutoff;
finite-cutoff admissibility at $Z=26$ is not a cutoff-uniform
critical-charge assertion.

\subsection{Positivity methods and what they actually supply}
Miyao~\cite{Miyao}, Theorem~3.4, proves finite-source Gaussian
domination for a half-filled Hubbard model on a bipartite lattice,
with two spin species, positive-type density interactions, and photons
entering through Peierls phases. Its particle--hole conjugate-factor
structure is not the spin-mixing projected Dirac operator here.
The thermodynamic susceptibility result also has an interaction
summability hypothesis which unscreened Coulomb interactions do not
satisfy. Our finite-CAR positivity proof and canonical projection are
therefore supplied explicitly rather than imported from that theorem.

Gaussian domination for periodic classical ferromagnets appears in
Fr\"ohlich--Simon--Spencer~\cite{FSS}. Euclidean vector-like
determinant positivity follows from paired anti-Hermitian eigenvalues
with a positive scalar mass, as in Vafa--Witten~\cite{VafaWitten}.
The Majorana positivity criteria of Wei et al.~\cite{Wei} and
Li--Jiang--Yao~\cite{LiJiangYao} provide further sufficient
auxiliary-field sign conditions. They are relevant precedents,
not substitutes for checking the present noncommuting projected
densities, spatial mass coefficient, charge twist, and photon integration.
No equivalence to an unverified Majorana criterion or priority claim
for the positivity mechanism is needed.

The componentwise Hubbard--Stratonovich formula is justified by the
spectral theorem for one Hermitian charge at a time. Trotter ordering
then handles their noncommutativity. This is not an operator-valued
Graf--Schenker localization theorem. It supplies finite-source
domination rather than only the curvature at the unsourced vacuum.
Indeed, $f(t)=\inf\sigma(H_0+tW)$ is concave, and its drop to
$t=1$ depends on curvature and possible level crossings along the
whole path. For $H_0=\operatorname{diag}(0,\Delta)$ and
$W=\operatorname{diag}(0,-B)$ with $B>\Delta>0$, the response
at zero vanishes while $f(1)=\Delta-B<0$.

Likewise an equal-time structure factor and a static susceptibility
are distinct quantities. With their connected spectral conventions,
\[
 S(k)=V^{-1}\|\rho(k)\Omega\|^2,\quad
 \chi(k)=V^{-1}\langle\rho(k)\Omega,(H-E)^{-1}_{\perp}
                         \rho(k)\Omega\rangle,
\]
Cauchy--Schwarz gives $S(k)^2\le\chi(k)F(k)$, where
$F(k)=V^{-1}\langle\rho(k)\Omega,(H-E)\rho(k)\Omega\rangle$.
A susceptibility estimate alone is not an equal-time bound; a lower
gap yields an inequality in the opposite direction. No such unproved
screening inference is used above.

\subsection{What remains outside these theorems}
The fixed-cutoff source bound grows as $\Lambda_f$ and the source-free
bare model with massless photons has the obstruction
\eqref{eq:vacuum-instability}. The mean-field point-source construction
proves necessary linear cutoff growth in its own reference convention,
not in every renormalized model or after the exact many-body vacuum
subtraction. The positive bare $A^2$ theorem and the exchange calculation
are likewise statements about the explicitly modified Hamiltonian.
The Pauli--Villars classical-field construction~\cite{GHLS2013} and
reduced-BDF charge renormalization~\cite{GLS2009} are different
frameworks and do not construct the present quantized theory's
ultraviolet limit.

A charge gap is not a neutral pair gap. A low-energy Pauli--Fierz or
Coulomb reduction would need a specified energy window, control of
the full pair-block resolvent and induced operators, and errors that
remain controlled with system size after vacuum subtraction. The
stability of a proposed effective Hamiltonian does not supply these
missing comparisons. Atomic binding, molecular binding, thermodynamic
limits of relative source energies, and equilibrium crystals each
require further arguments. The purpose of this paper is to preserve
the proved regulator-dependent comparisons and their counterexamples,
without identifying them with those additional conclusions.
